\subsubsection{専用言語と領域固有言語}

設計表現は、構造と振る舞いの二分法だけで整理できるとは限らない。たとえば、利用者インタフェース設計では、画面の構造、操作の流れ、入力制約、表示規則が混ざる。また、安全性や信頼性の領域では、その分野に特化した記法が発達している。

\term{領域固有言語}{Domain-Specific Language, DSL} は、特定の業務領域や技術領域の概念を直接表すための言語である。汎用プログラミング言語ではなく、その領域の語彙を使って設計や実装に近い表現を行う。例として、シミュレーション、リアルタイムシステム、ゲーム開発、利用者インタフェース、テスト開発、言語処理などの領域向けDSLがある。

DSLの利点は、領域の専門家が理解しやすい語彙で表現できることである。一方、DSLを作るには、言語仕様、構文解析、生成、保守の負担がある。したがって、繰り返し使う価値がある領域で特に有効である。
